% Part: computability
% Chapter: recursive-functions
% Section: pr-relations

\documentclass[../../../include/open-logic-section]{subfiles}

\begin{document}

\olfileid{cmp}{rec}{prr}
\olsection{原始递归关系}


\begin{defn}
称关系 $R(\vec x)$ 是\emph{原始递归}的，如果其特征函数
\[
\Char{R}(\vec x) = \left\{
  \begin{array}{ll}
  1 & \mbox{若 $R(\vec x)$} \\
  0 & \mbox{否则}
  \end{array}
\right.
\]
是原始递归的。
\end{defn}

换言之，所谓原始递归关系 $R(\vec x)$，指的是形如 $\Char{R}(\vec x) = 1$ 的关系，其中 $\Char{R}$ 是一个在任何输入下均返回 1 或 0 的原始递归函数。例如，关系 $\fn{IsZero}(x)$ 成立当且仅当 $x = 0$，它对应于函数 $\Char{\fn{IsZero}}$，该函数由原始递归定义为
\begin{align*}
\Char{\fn{IsZero}}(0) & = 1,\\
\Char{\fn{IsZero}}(x+1) & = 0.
\end{align*}

显然，关系可以与其他原始递归函数复合。因此，以下关系也是原始递归的：
\begin{enumerate}
\item 相等关系 $x = y$，定义为 $\fn{IsZero}(\left|x -
  y\right|)$
\item 小于等于关系 $x \leq y$，定义为 $\fn{IsZero}(x \tsub y)$
\end{enumerate}


\begin{prop}
  原始递归关系的集合对布尔运算封闭，即，
  若 $P(\vec x)$ 和 $Q(\vec x)$ 是原始递归的，则以下也是原始递归的：
  \begin{enumerate}
  \item $\lnot P(\vec x)$
  \item $P(\vec x) \land Q(\vec x)$
  \item $P(\vec x) \lor Q(\vec x)$
  \item $P(\vec x) \lif Q(\vec x)$
  \end{enumerate}
\end{prop}

\begin{proof}
  假设 $P(\vec x)$ 和 $Q(\vec x)$ 是原始递归的，即它们的特征函数 $\Char{P}$ 和 $\Char{Q}$ 是原始递归的。我们只需证明 $\lnot P(\vec x)$ 等关系的特征函数也是原始递归的。
  \[
  \Char{\lnot P}(\vec x) = \begin{cases}
    0 & \text{若 $\Char{P}(\vec x) = 1$}\\
    1 & \text{否则}
  \end{cases}
  \]
  我们可以将 $\Char{\lnot P}(\vec x)$ 定义为 $1 \tsub \Char{P}(\vec x)$。
  \[
  \Char{P \land Q}(\vec x) = \begin{cases}
    1 & \text{若 $\Char{P}(\vec x) = \Char{Q}(\vec x) = 1$}\\
    0 & \text{否则}
  \end{cases}
  \]
  我们可以将 $\Char{P \land Q}(\vec x)$ 定义为 $\Char{P}(\vec x) \cdot
  \Char{Q}(\vec x)$，或定义为 $\fn{min}(\Char{P}(\vec x), \Char{Q}(\vec
  x))$。类似地，
  \begin{align*}
    \Char{P \lor Q}(\vec x) & = \fn{max}(\Char{P}(\vec x), \Char{Q}(\vec x)) \text{ 与}\\
    \Char{P \lif Q}(\vec x) & = \fn{max}(1 \tsub
  \Char{P}(\vec x), \Char{Q}(\vec x)).
  \end{align*}
\end{proof}

\begin{prop}
  原始递归关系的集合对有界量词封闭，即如果 $R(\vec x, z)$ 是原始递归关系，那么关系
  \begin{align*}
    & \bforall{z < y}{R(\vec x, z)} \text{ 与}\\
    & \bexists{z < y}{R(\vec x, z)}.
  \end{align*}
  也是原始递归的。
  $\bforall{z < y}{R(\vec x, z)}$ 对 $\vec x$ 和 $y$ 成立当且仅当 $R(\vec x, z)$ 对每个小于 $y$ 的 $z$ 都成立，对 $\bexists{z < y}{R(\vec x, z)}$ 亦然。
\end{prop}

\begin{proof}
  按约定，我们取 $\bforall{z < 0}{R(\vec x, z)}$ 为真（出于不存在小于 $0$ 的 $z$ 这一平凡理由），而取 $\bexists{z < 0}{R(\vec x, z)}$ 为假。有界全称量词的作用相当于有限积或迭代最小值，即如果 $P(\vec x, y) \defiff \bforall{z <
  y}{R(\vec x, z)}$，那么 $\Char{P}(\vec x, y)$ 可由下式定义：
  \begin{align*}
    \Char{P}(\vec x, 0) & = 1\\
    \Char{P}(\vec x, y+1) & =
    \fn{min}(\Char{P}(\vec x, y), \Char{R}(\vec x, y)).
  \end{align*}
  有界存在量词可以类似地用 $\fn{max}$ 来定义。或者，也可以利用等价式 $\bexists{z < y}{R(\vec x, z)}
  \liff \lnot \bforall{z < y}{\lnot R(\vec x, z)}$ 通过有界全称量词来定义。注意，例如，形式为 $\bexists{x \leq y}{\dots
  x\dots}$ 的有界量词等价于 $\bexists{x < y+1}{\dots x \dots}$。
\end{proof}

\begin{prob}
  证明三元关系 $x \equiv y \mod n$（模 $n$ 同余）是原始递归的。
\end{prob}

另一个有用的原始递归函数是条件函数 $\fn{cond}(x,y,z)$，其定义为
\begin{align*}
  \fn{cond}(x,y,z) & = \begin{cases}
  y & \text{若 $x = 0$} \\
  z & \text{否则}.
\end{cases}
\intertext{其递归定义为}
\fn{cond}(0,y,z) & = y,\\
\fn{cond}(x+1,y,z) & = z.
\end{align*}
我们可以利用它来论证由原始递归关系分情况定义的函数仍是原始递归的：

\begin{prop}
如果 $g_0(\vec x)$, \dots,~$g_m(\vec x)$ 是原始递归函数，而 $R_0(\vec x)$, \dots, $R_{m-1}(\vec x)$ 是原始递归关系，那么由下式定义的函数 $f$
\[
f(\vec x) = \begin{cases}
    g_0(\vec x) & \text{若 $R_0(\vec{x})$} \\
    g_1(\vec x) & \text{若 $R_1(\vec{x})$ 且非 $R_0(\vec{x})$} \\
    \vdots & \\
    g_{m-1}(\vec x) & \text{若 $R_{m-1}(\vec{x})$ 且前面的条件均不满足}
    \\
    g_m(\vec x) & \mbox{否则}
\end{cases}
\]
也是原始递归的。
\end{prop}

\begin{proof}
  当 $m = 1$ 时，这就是由 \[
  f(\vec x) = \fn{cond}(\Char{\lnot R_0}(\vec x),g_0(\vec x),g_1(\vec
  x)).
  \] 定义的函数。
  当 $m > 1$ 时，只需复合这种形式的定义即可。
\end{proof}
\end{document}